\documentclass[a4paper,12pt]{article}
\usepackage[T2A]{fontenc}
\usepackage[utf8]{inputenc}
\usepackage[russian]{babel}
\usepackage{amsmath,amssymb}
\usepackage{helvet}
\renewcommand{\familydefault}{\sfdefault}
\usepackage{geometry}
\geometry{left=22mm,right=22mm,top=22mm,bottom=22mm}
\usepackage{microtype}
\usepackage{tabularx,booktabs,array}
\usepackage{enumitem}
\usepackage{fancyhdr}
\usepackage{parskip}
\usepackage{xcolor}
\usepackage{tikz}
\usepackage{hyperref}
\hypersetup{hidelinks}
\definecolor{accent}{HTML}{008080}
\definecolor{deepaccent}{HTML}{005D5D}
\definecolor{textgray}{HTML}{353C3C}
\definecolor{pale}{HTML}{E8F3F1}
\color{textgray}
\setlength{\parindent}{0pt}
\setlength{\parskip}{5pt}
\setlength{\headheight}{15pt}
\pagestyle{fancy}
\fancyhf{}
\fancyhead[L]{\footnotesize МАТЕМАТИКА\quad\textcolor{accent}{|}\quad Ярослав Верещагин}
\fancyhead[R]{\footnotesize 12.11.2026}
\fancyfoot[L]{\footnotesize\textcolor{accent}{Лёвин А.~А.}}
\fancyfoot[R]{\footnotesize\thepage}
\renewcommand{\headrulewidth}{0pt}
\renewcommand{\footrulewidth}{0pt}
\newcommand{\Cover}[3]{%
  \thispagestyle{empty}
  \begin{center}
  \vspace*{20mm}
  {\large\color{deepaccent} МАТЕМАТИКА\par}
  \vspace{21mm}
  {\LARGE #1\par}
  \vspace{11mm}
  {\Large\color{accent} #2\par}
  \vspace{20mm}
  {\large Подготовка к ОГЭ по математике\par}
  \vspace{40mm}
  {\large Ярослав Верещагин\par}
  \vspace{15mm}
  {\large 12 ноября 2026 г.\par}
  \vfill
  {\small Лёвин Артём Александрович\par}
  {\small эксклюзивно для Ярослава Верещагина\par}
  \vspace{15mm}
  \begin{tikzpicture}
     \draw[color=accent, line width=.75pt] (0,0) .. controls (1.3,.55) and (2.8,-.42) .. (4.0,0)
      .. controls (5.0,.47) and (6.1,-.35) .. (7.1,0)
      .. controls (8.1,.41) and (9.1,-.48) .. (10.5,0);
  \end{tikzpicture}
  \end{center}
  \clearpage
}
\newcommand{\Section}[1]{\vspace{2mm}{\large\bfseries\color{deepaccent}#1}\par\vspace{2mm}}
\newcommand{\Closing}{\vspace{8mm}\begin{center}{\footnotesize\color{deepaccent}Лёвин Артём Александрович\par эксклюзивно для Ярослава Верещагина}\end{center}}
\begin{document}

\Cover{Таблица сверки ответов}{Вариант ОГЭ\\[4pt]\large Фиксация ответов и границы проверки}{}
\Section{Ответы ученика}
\textbf{Основание:} фотография листа ответов, приложенная в папке с работами. В ней приведены ответы к номерам 1--19.

\textbf{Важно.} В папке файл с названием «Вариант ОГЭ.pdf» содержит чек-лист «Движение по реке» от 9 ноября 2026 г. с \textbf{десятью} упражнениями. Исходные условия варианта ОГЭ с номерами 1--19 и официальный ключ к ним отсутствуют. Поэтому окончательно оценить правильность ответов и достоверно определить темы всех заданий пока невозможно.

\small
\renewcommand{\arraystretch}{1.28}
\begin{tabularx}{\linewidth}{@{}>{\centering\arraybackslash}p{11mm} >{\raggedright\arraybackslash}p{46mm} >{\centering\arraybackslash}p{36mm} X@{}}
\toprule
\textbf{№} & \textbf{Ответ ученика} & \textbf{Правильный ответ} & \textbf{Тема задания} \\
\midrule
1 & Не до конца разборчиво & --- & Не установлена \\
2 & $7{,}92$ & --- & Десятичные дроби* \\
3 & $2\,000$ & --- & Арифметические действия* \\
4 & $16\,750$ руб. & --- & Не установлена \\
5 & 33 & --- & Не установлена \\
6 & 4 & --- & Не установлена \\
7 & 1 & --- & Не установлена \\
8 & $13{,}5\sqrt{2}$ & --- & Корни и выражения* \\
9 & $-\dfrac{36}{23}$ & --- & Обыкновенные дроби* \\
10 & --- & --- & Не установлена \\
11 & 321 & --- & Не установлена \\
12 & 304 & --- & Не установлена \\
13 & 2 & --- & Не установлена \\
14 & --- & --- & Не установлена \\
15 & --- & --- & Не установлена \\
16 & --- & --- & Не установлена \\
17 & 32 & --- & Не установлена \\
18 & 5 & --- & Не установлена \\
19 & 123 & --- & Не установлена \\
\bottomrule
\end{tabularx}
\normalsize

\vspace{3mm}
\footnotesize
\textbf{Условные обозначения:} «---» в графе правильного ответа означает, что ответ невозможно установить без условия задания. «---» в графе ответа ученика означает отсутствие записи на фотографии. Звёздочка у темы означает, что она предположена \emph{только по характеру записи}; тема исходного экзаменационного задания неизвестна.
\normalsize

\Section{Отдельная проверка арифметики}
В записи к \textbf{№ 3} указано равенство \(26\,500-19\,000+5\,500=2\,000\). Оно \textbf{арифметически неверно}: по порядку действий левая часть равна \(13\,000\). Это подтверждённая ошибка вычисления, однако окончательный ответ на \emph{исходную} задачу № 3 без её условия определить нельзя.

\textbf{Для завершения сверки} необходимы страницы исходного варианта ОГЭ или официальный ключ с условием каждого задания. По № 1 желательно также уточнить написанный ответ.
\end{document}
